\documentclass[11pt]{article}

\usepackage[T1]{fontenc}
\usepackage[utf8]{inputenc}
\usepackage{amsmath,amssymb,amsthm}
\usepackage{booktabs}
\usepackage{geometry}
\usepackage{hyperref}
\usepackage{microtype}
\usepackage{xcolor}

\geometry{margin=1in}
\setlength{\emergencystretch}{2em}
\hypersetup{
  colorlinks=true,
  linkcolor=blue!55!black,
  citecolor=blue!55!black,
  urlcolor=blue!55!black
}

\newtheorem{theorem}{Theorem}[section]
\newtheorem{lemma}[theorem]{Lemma}
\newtheorem{proposition}[theorem]{Proposition}
\newtheorem{corollary}[theorem]{Corollary}
\theoremstyle{definition}
\newtheorem{definition}[theorem]{Definition}
\theoremstyle{remark}
\newtheorem{remark}[theorem]{Remark}

\newcommand{\F}{\mathbb F}
\newcommand{\A}{\mathcal A}
\newcommand{\Q}{\mathcal Q}
\newcommand{\C}{\mathcal C}
\newcommand{\Span}{\operatorname{span}}
\newcommand{\FiveCDC}{\textup{FiveCDC}}
\newcommand{\one}{\mathbf 1}
\numberwithin{equation}{section}

\title{Parallel Six-Point Compression and a\\
Component--Cycle Span Criterion for\\
Five-Cycle Double Covers}
\author{Atharva Vaidya\\
\small AI-assisted research draft for independent human verification}
\date{July 29, 2026}

\begin{document}
\maketitle

\begin{center}
\fcolorbox{red!70!black}{red!4}{
\parbox{0.91\textwidth}{
\textbf{Resolution status, trust boundary, and AI use.}
The standard five-cycle double cover conjecture remains open.  This
manuscript proves an exact reformulation for loopless cubic multigraphs:
after a nowhere-zero \(\F_2^3\)-flow, a Fano plane, and a direction have
been chosen, a parallel six-point compression is soluble exactly when an
explicit component defect belongs to a circuit-switch span.  Existentially
choosing those data is equivalent to \FiveCDC; no universal choice is
proved.  The orientable conjecture is not considered.

OpenAI Codex agents using GPT-5-series models, under Atharva Vaidya's
direction, discovered the formulation, derived the proofs, wrote the
producer and checking programs, ran the computations, conducted an
AI-assisted prior-art comparison, and drafted this manuscript.  The
checker is structurally independent of the discovery implementation, but
it is also AI-authored.  No human peer review or priority determination is
claimed.  A human author must check every argument, rerun or independently
reimplement the audits, complete a specialist literature review, and
accept scholarly responsibility before submission.
}}
\end{center}

\begin{abstract}
Let \(W=\F_2^3\), let \(G\) be a finite connected loopless cubic
multigraph, and fix a nowhere-zero \(W\)-flow \(f\).  A companion
six-point one-hole construction asks for a coordinate-pair labelling
supported on six points of \(W\), avoiding one additional pair which is
then merged.  We give a purely combinatorial characterization for the
21 translation classes in which the omitted pair and merge pair are
parallel.

Fix a two-plane \(H\le W\) and \(0\ne d\in H\).  The edges with values in
\(H\) form components \(Q\) of degree one or three; the complementary
edges form circuits on their leaves.  A cut defect
\(\beta_H\in\F_2^{\Q}\) is determined by any one outside value.  Each
outside circuit \(C\) supplies a switch vector
\(\gamma_C^d\), recording the components in which \(C\) has an odd
number of leaves whose inside-edge value differs from \(d\).  We prove
that the parallel six-point system is soluble if and only if
\[
                  \beta_H\in\Span_{\F_2}\{\gamma_C^d:C\}.
\]
These columns are literal flow moves.  Adding \(d\) to every edge of a
chosen outside circuit preserves a nowhere-zero flow and leaves its
\(H\)-edge subgraph fixed; it changes \(\beta_H\) by exactly
\(\gamma_C^d\).  Thus the parallel system is soluble exactly when a set
of whole-circuit \(d\)-switches reaches the Hu\v{s}ek--\v{S}\'amal
five-point condition \(\beta_H=0\).  Equivalently, on every outside
circuit one chooses one of two \(d\)-paired colour classes so that the
chosen edges become Eulerian after the \(H\)-components are contracted.

We also prove that existentially choosing \(f,H,d\) in this criterion is
equivalent to the existence of a standard five-cycle double cover.  We
explain precisely how this differs from Hu\v{s}ek--\v{S}\'amal,
Theorem~3.16: their five-point condition is the special case
\(\beta_H=0\), whereas the circuit columns can cancel a nonzero defect.
The graph-level existential statements are nevertheless equivalent.

A standard-library checker independently generates all 900
\(\mathrm{GL}(3,2)\)-orbits of nowhere-zero flows on a rigid
12-vertex graph and compares the original affine and reduced decisions
on all 18,900 parallel systems, with no disagreement.  A separate
AI-authored SAT census records an odd 21-flag selector in every feasible
star fibre of all 4,461 simple biconnected cubic graphs through order 16.
These computations audit and motivate the reformulation; they do not
resolve the universal conjecture.
\end{abstract}

\section{Definitions, conventions, and exact scope}

A \emph{cycle} in this paper is an Eulerian edge-subset, not necessarily
connected or \(2\)-regular.  A labelled five-cycle double cover is an
indexed family \((C_1,\ldots,C_5)\) of cycles such that every edge lies
in exactly two members.  Empty members are allowed, so ``five'' means
``at most five.''  This agrees with the formulation of
Hu\v{s}ek--\v{S}\'amal~\cite[Definition~1.2]{HusekSamal2026}.
We work with finite connected loopless cubic multigraphs.  Parallel
edges are distinct; loops are excluded.  General-graph reductions and
the orientable variant are outside the present statement.

Put \(W=\F_2^3\), written additively.  A map
\[
                      f:E(G)\longrightarrow W-\{0\}
\]
is a nowhere-zero \(W\)-flow when
\[
                      \sum_{e\ni v}f(e)=0
                      \qquad(v\in V(G)).
\]
At a cubic vertex the three incident values are distinct and form the
nonzero part of a unique two-plane \(L_v\le W\).

The companion preprint~\cite{VaidyaSixPoint2026} proves the following
potential description.  A compatible pair labelling assigns
\[
                       P_e\in\binom W2,\qquad
                       \sum_{x\in P_e}x=f(e),
\]
such that each coordinate has even local degree.  For an incidence
\(e\ni v\), put \(p=f(e)\) and choose either of the other two incident
flow values \(s_{v,e}\).  All compatible labellings are parameterized
by potentials \(t_v\in W\), with
\begin{equation}
 P_{v,e}=t_v+s_{v,e}+\langle p\rangle,                       \label{eq:side}
\end{equation}
subject to
\begin{equation}
 t_u+t_v+s_{u,e}+s_{v,e}\in\langle f(e)\rangle
 \qquad(e=uv).                                               \label{eq:edgecompat}
\end{equation}
The choice of \(s_{v,e}\) does not affect the pair.

Choose disjoint pairs \(R,M\in\binom W2\).  The six-point one-hole
system additionally requires every local triangle to avoid the points
of \(R\) and every edge label to differ from \(M\).  Identifying the two
members of the unused pair \(M\) then gives five coordinate names
\cite{VaidyaSixPoint2026}.  Translation of all eight points preserves
the problem.

We study the \emph{parallel} classes.  After translation they have the
unique normal form
\begin{equation}
 R=\langle d\rangle=\{0,d\},\qquad
 M=a+\langle d\rangle=\{a,a+d\},\qquad
 H=\langle a,d\rangle,                                      \label{eq:normal}
\end{equation}
where \(H\le W\) is a two-plane and \(0\ne d\in H\).  There are seven
planes and three nonzero directions in each, hence 21 parallel flags
\((d,H)\).

\begin{remark}[provisional novelty statement]
Hu\v{s}ek--\v{S}\'amal already prove an exact flow formulation of
\FiveCDC and the component criterion discussed in
Section~\ref{sec:HS}.  The companion preprint already proves the general
six-point affine system and its parallel/skew local tables.  The possible
additional contribution here is narrower: the component--circuit
elimination of the \emph{parallel} six-point system, the exact
whole-circuit switch normal form relating it back to Theorem~3.16, and
the elementary proof that the parallel subclass alone remains graph-level
exact.  We have not found this precise factorization in the cited sources,
but our search is incomplete and we make no claim of priority.
\end{remark}

\section{The component and circuit data}

Fix \(f,H,d\) as above and set
\begin{equation}
 E_H=\{e:f(e)\in H\},\qquad F=E(G)-E_H.                     \label{eq:EH}
\end{equation}
At a vertex \(v\), either \(L_v=H\), in which case its \(E_H\)-degree is
three, or \(L_v\cap H\) is a line, in which case its \(E_H\)-degree is
one.  Call the latter vertices \emph{leaves}.  The graph \(F\) has
degree two at every leaf and degree zero elsewhere, so its nonempty
components are circuits.  Write \(\Q\) for the components of
\((V(G),E_H)\) and \(\C\) for the \(F\)-circuits.

For a leaf \(v\), let \(p_v\) be the value of its unique incident
\(E_H\)-edge.  For \(Q\in\Q\) and \(C\in\C\), define
\begin{equation}
 \gamma_C^d(Q)=
 \left|\{v\in V(Q)\cap V(C):p_v\ne d\}\right|\pmod2.         \label{eq:gamma}
\end{equation}

Choose any \(w\in W-H\), and define
\begin{equation}
 \beta_H(Q)=
 \left|\delta_G(Q)\cap F\cap f^{-1}(w)\right|\pmod2.         \label{eq:beta}
\end{equation}
The next lemma makes this well defined.

\begin{lemma}[outside-value independence]\label{lem:beta}
For every \(Q\in\Q\), the value in \eqref{eq:beta} is independent of
the choice \(w\in W-H\).
\end{lemma}

\begin{proof}
Every vertex has degree one or three in \(Q\).  The handshaking lemma
therefore gives \(|V(Q)|\equiv0\pmod2\).  Since \(G\) is cubic,
\[
 |\delta_G(Q)|\equiv 3|V(Q)|\equiv0\pmod2.
\]
Every edge of \(\delta_G(Q)\) lies in \(F\), hence has one of the four
values outside \(H\).  Write \(H=\{0,p,q,p+q\}\) and choose
\(w_0\notin H\).  Let \(n_0,n_p,n_q,n_{p+q}\) be the four boundary
multiplicities modulo two.  Their sum is zero by the preceding
even-cardinality equation.  Summing the flow equations over the
vertices of \(Q\) gives
\[
 n_0w_0+n_p(w_0+p)+n_q(w_0+q)+n_{p+q}(w_0+p+q)=0.
\]
Comparing coefficients in the basis \((w_0,p,q)\) shows
\[
 n_p=n_q=n_{p+q}=n_0.
\]
Thus all four outside colour classes have the same boundary parity.
\end{proof}

\section{The parallel component criterion}

The local lists for the parallel normal form
\eqref{eq:normal}, proved in the companion paper, are
\begin{equation}
\A_L(R,M)=
\begin{cases}
W-H,&L=H,\\
M,&d\in L,\ L\ne H,\\
R,&d\notin L.
\end{cases}                                                  \label{eq:locallist}
\end{equation}
For completeness, this table can be checked directly.  A local
triangle \((t+L)-\{t\}\) must avoid \(R\) and must not have \(M\) as a
side.  If \(d\in L\), the two omitted points lie in one \(L\)-coset; if
\(d\notin L\), they lie in opposite cosets.  Separating \(L=H\) from
the other cases gives \eqref{eq:locallist}.

\begin{theorem}[parallel component criterion]\label{thm:component}
For a fixed nowhere-zero \(W\)-flow \(f\) and parallel flag \((d,H)\),
the six-point one-hole potential system
\eqref{eq:side}--\eqref{eq:locallist} is soluble if and only if
\begin{equation}
             \beta_H\in
             \Span_{\F_2}\{\gamma_C^d:C\in\C\}
             \subseteq\F_2^{\Q}.                            \label{eq:span}
\end{equation}
\end{theorem}

\begin{proof}
Choose \(b\notin H\).  Let \(v\) be a leaf, with inside-edge value
\(p_v\).  By \eqref{eq:locallist},
\[
                         y_v:=t_v+p_v\in M.
\]
Thus there is a unique \(z_v\in\F_2\) such that
\begin{equation}
                         y_v=a+z_vd.                         \label{eq:y}
\end{equation}

Consider an \(F\)-edge of value \(q\) joining leaves \(u,v\).  At
either endpoint, formula \eqref{eq:side}, using \(p_u\) or \(p_v\) as
the other incident value, writes its side as
\(y+\langle q\rangle\).  The endpoint sides agree exactly when
\[
                         y_u+y_v\in\langle q\rangle.
\]
The left side belongs to \(H\), while \(q\notin H\), so this holds
exactly when \(y_u=y_v\).  Hence \(z_v\) is a single common variable
\(z_C\) on each circuit \(C\in\C\).

We now encode sides on the \(E_H\)-edges.  An outside pair of difference
\(p\in H-\{0\}\) is one of the two cosets of
\(\langle p\rangle\) in \(b+H\).  Give
\(b+\langle p\rangle\) port bit one and the other pair port bit zero.
At a degree-three \(E_H\)-vertex, the four potentials \(t\in b+H\)
give precisely the four even-parity triples of incident port bits.
Equivalently, writing \(t=b+h\) and the incident values as
\(p,q,p+q\), a direct substitution in \eqref{eq:side} shows that the
three port bits have xor zero.

At a leaf \(v\), write either incident outside value as \(q=b+h\).
The port bit on the \(p_v\)-edge equals one precisely when
\begin{equation}
                     a+h+z_Cd\in\langle p_v\rangle.          \label{eq:leafport}
\end{equation}
At \(z_C=0\), condition \eqref{eq:leafport} holds exactly when one of
the two incident \(F\)-edge values is \(b+a\).  Adding \(d\) toggles
the condition exactly when \(p_v\ne d\).  Consequently the forced
leaf bit is
\begin{equation}
 \ell_v(z_C)=c_v+z_C\one[p_v\ne d],                          \label{eq:ell}
\end{equation}
where \(c_v=1\) precisely when an incident outside edge has value
\(b+a\).

Within a connected component \(Q\) of \(E_H\), common edge port bits
satisfying xor zero at every degree-three vertex exist exactly when
\begin{equation}
                  \sum_{\substack{v\in V(Q)\\v\ {\rm leaf}}}
                  \ell_v(z_{C(v)})=0.                        \label{eq:component}
\end{equation}
Necessity follows by summing all vertex equations.  For sufficiency,
solve edge bits successively along a spanning tree of \(Q\); the only
remaining equation is \eqref{eq:component}.  A single edge between two
leaves has the same equality condition, and a component without leaves
has no obstruction.

The coefficient of \(z_C\) in \eqref{eq:component} is
\(\gamma_C^d(Q)\).  Its constant term counts incidences in \(Q\) of
\((b+a)\)-valued \(F\)-edges.  An edge internal to \(Q\) contributes
twice and a boundary edge once, so the constant term is
\(\beta_H(Q)\), by Lemma~\ref{lem:beta}.  Thus all component equations
are the binary system
\[
                    \sum_{C\in\C}\gamma_C^d z_C=\beta_H.
\]
It is soluble exactly when \eqref{eq:span} holds.
\end{proof}

\section{The circuit-switch normal form}

Contract every component \(Q\in\Q\) to one vertex, retaining the
resulting loops and parallel edges.  For \(C\in\C\) and \(x\notin H\),
let
\[
                   \delta_C^x\in\F_2^{\Q}
\]
be the odd-degree boundary vector of the \(x\)-valued edges of \(C\) in
this quotient.

\begin{theorem}[whole-circuit switch normal form]\label{thm:switch}
For \(z=(z_C:C\in\C)\in\F_2^\C\), define
\begin{equation}
 f^z(e)=
 \begin{cases}
 f(e)+z_Cd,&e\in E(C),\ C\in\C,\\
 f(e),&e\in E_H.
 \end{cases}                                                 \label{eq:switch}
\end{equation}
Then \(f^z\) is a nowhere-zero \(W\)-flow,
\(E_H(f^z)=E_H(f)\), and
\begin{equation}
             \beta_H(f^z)=
             \beta_H(f)+\sum_{C\in\C}z_C\gamma_C^d.          \label{eq:betaswitch}
\end{equation}
Consequently the parallel six-point system for \((f,d,H)\) is soluble
if and only if a set of whole outside-circuit \(d\)-switches produces a
flow \(f^z\) satisfying the Hu\v{s}ek--\v{S}\'amal five-point condition
\(\beta_H(f^z)=0\).
\end{theorem}

\begin{proof}
Every leaf of \(E_H\) is incident with two edges of its unique
\(F\)-circuit, so a selected circuit adds \(d+d=0\) to the flow sum
there.  It changes no edge at a degree-three \(E_H\)-vertex.  Therefore
\(f^z\) is a flow.  A value outside \(H\), when translated by
\(d\in H\), remains outside \(H\) and in particular remains nonzero.
Thus \(f^z\) is nowhere zero and its \(H\)-valued edge set is unchanged.

Fix \(w\notin H\).  On a selected circuit, adding \(d\) swaps the
\(w\)- and \((w+d)\)-valued edge classes.  The change of the defect is
therefore \(\delta_C^w+\delta_C^{w+d}\).  At a leaf \(v\in C\), its two
outside values differ by \(p_v\); they meet the \(d\)-pair
\(\{w,w+d\}\) in odd parity exactly when \(p_v\ne d\).  After internal
edges of each \(Q\) cancel twice, this boundary difference is
\(\gamma_C^d\).  Summing over selected circuits proves
\eqref{eq:betaswitch}.  Finally,
\[
 \beta_H(f^z)=0
 \quad\Longleftrightarrow\quad
 \beta_H(f)\in\Span\{\gamma_C^d:C\in\C\},
\]
so Theorem~\ref{thm:component} proves the last assertion.
\end{proof}

\begin{proposition}[circuit-colour dictionary]\label{prop:dictionary}
For every \(C\in\C\), every \(d\in H-\{0\}\), and every \(w\notin H\),
\begin{equation}
 \gamma_C^d=\delta_C^w+\delta_C^{w+d},
 \qquad
 \beta_H=\sum_{C\in\C}\delta_C^w.                            \label{eq:dictionary}
\end{equation}
Consequently, the parallel system for \((d,H)\) is soluble if and only
if one can choose independently on each circuit \(C\) either the
\(w\)-valued or the \((w+d)\)-valued edges so that the union of all
chosen edges is Eulerian in the contracted quotient.
\end{proposition}

\begin{proof}
At a leaf \(v\in C\), the two incident outside values differ by
\(p_v\).  Their intersection with the \(d\)-pair
\(\{w,w+d\}\) has even size if \(p_v=d\) and odd size otherwise.
Summing this incidence parity over the leaves of \(C\) in each \(Q\),
with internal edges counted twice, proves the first identity.  The
second identity is just \eqref{eq:beta} split circuit by circuit.

Now \(\beta_H\) belongs to the span of the \(\gamma_C^d\) exactly when
there are \(z_C\in\F_2\) with
\[
 0=\beta_H+\sum_C z_C\gamma_C^d
  =\sum_C\bigl((1+z_C)\delta_C^w+
                     z_C\delta_C^{w+d}\bigr).
\]
This says precisely that the selected union has zero boundary vector.
\end{proof}

\section{The parallel normal form is graph-level exact}

\begin{theorem}[exact parallel reformulation]\label{thm:exact}
A finite connected loopless cubic multigraph \(G\) has a standard
five-cycle double cover if and only if there are a nowhere-zero
\(W\)-flow \(f\), a two-plane \(H\le W\), and
\(0\ne d\in H\) satisfying \eqref{eq:span}.
\end{theorem}

\begin{proof}
A five-cycle double cover is equivalently an edge labelling by
two-subsets of five coordinate names such that each name has even local
degree.  At a cubic vertex the three duads must be the three sides of a
triangle on three names: parity makes their symmetric difference empty,
and no two can coincide.  Conversely such duad labels recover the five
Eulerian coordinate edge-sets.

Suppose first that \(G\) has such a labelling.  Inject the five coordinate
names into \(W\), leaving three unused points
\[
                              B=\{u,v,w\}.
\]
Give an edge labelled \(\{x,y\}\) the value \(x+y\).  The local triangle
condition makes this a nowhere-zero \(W\)-flow.  Put
\begin{equation}
             R=\{u,v\},\qquad d=u+v,\qquad
             M=\{w,w+d\}.                                  \label{eq:unused}
\end{equation}
These pairs are parallel and disjoint.  Indeed, \(w\ne u,v\);
\(w+d=u\) would imply \(w=v\), and \(w+d=v\) would imply \(w=u\).
Every original triangle avoids \(R\), and no original edge label equals
\(M\), because \(u,v,w\) are unused.  Thus the original labelling is a
solution of this parallel six-point system.  Translating by \(u\) puts
\eqref{eq:unused} in the normal form \eqref{eq:normal}, with
\[
                             H=\langle d,w+u\rangle.
\]
Theorem~\ref{thm:component} gives \eqref{eq:span}.

Conversely, Theorems~\ref{thm:component} and~\ref{thm:switch} give a
whole-circuit switched flow \(f^z\) satisfying \(\beta_H(f^z)=0\).
By Hu\v{s}ek--\v{S}\'amal Theorem~3.16 and Observation~3.15, this
already gives a five-cycle double cover.

There is also a direct compression proof.  A solution of a parallel
system gives compatible duad labels on the six points \(W-R\), with
\(M\) unused.  Identify the two endpoints of \(M\).  No used duad
collapses, because \(M\) itself is unused.  Each unmerged coordinate
retains even local degree, and the merged coordinate has degree equal
modulo two to the sum of the two old degrees.  The quotient labels
therefore give five Eulerian edge-subsets covering every edge exactly
twice.
\end{proof}

\section{Exact relation to Hu\v{s}ek--\v{S}\'amal}
\label{sec:HS}

Hu\v{s}ek--\v{S}\'amal fix a nonzero functional
\(\mu\in W^*\), an outside value \(w\) with \(\mu(w)=1\), and set
\[
 F_{\rm HS}=\{e:\mu(f(e))=1\},\qquad
 M_{\rm HS}=f^{-1}(w),\qquad
 Z=\{v:v\hbox{ meets }M_{\rm HS}\}.
\]
With \(H=\ker\mu\), their graph
\((V,E\setminus F_{\rm HS})\) is exactly our \((V,E_H)\).
For a component \(Q\in\Q\),
\[
 |Z\cap V(Q)|
 \equiv|\delta_G(Q)\cap f^{-1}(w)|
 =\beta_H(Q)\pmod2,
\]
because an internal \(w\)-edge contributes two endpoints and a boundary
\(w\)-edge contributes one.

Their Theorem~3.16 states that their fixed-flow five-point system is
soluble exactly when every such component contains an even number of
vertices of \(Z\)~\cite[Theorem~3.16]{HusekSamal2026}.  In our notation,
this is
\begin{equation}
                              \beta_H=0.                     \label{eq:HS}
\end{equation}
It is the clean-defect special case of
Theorem~\ref{thm:component}: if \eqref{eq:HS} holds, then
\eqref{eq:span} holds for all three nonzero \(d\in H\).
For a fixed \(f,H\), our six-point system can also be soluble when
\(\beta_H(f)\ne0\), because the circuit columns can cancel the defect.
Theorem~\ref{thm:switch} gives the exact relationship: this is the
closure of the Hu\v{s}ek--\v{S}\'amal test under legal additions of
\(d\) on whole outside circuits.  The switched flow \(f^z\) has the
same \(E_H\)-components and satisfies \(\beta_H(f^z)=0\).  In the
coordinate-pair picture, the same freedom comes from allowing a sixth
coordinate and forbidding one merge pair rather than requiring three
coordinates to be empty.

The distinction must not be overstated.  Hu\v{s}ek--\v{S}\'amal already
prove that the existence of some flow and plane satisfying
\eqref{eq:HS} is equivalent to \FiveCDC
\cite[Observation~3.15 and Conjecture~3.19]{HusekSamal2026}.
Theorem~\ref{thm:exact} gives another existentially equivalent normal
form, not a stronger graph-level existence theorem.  Their
Theorem~3.16 also counts the solutions of the fixed-flow five-point
system; the present theorem records solubility, not an analogous count.

\section{Computational audits}
\label{sec:computation}

The proofs above are elementary and do not depend on computation.  The
following two experiments have different trust boundaries and are
reported separately.

\subsection{Independent fixed-flow regression audit}

The standard-library program \texttt{verify.py} in the accompanying
directory:
\begin{enumerate}
\item decodes the rigid 12-vertex cubic graph
      \texttt{K??FEaKR@oE\_} and checks connectedness and bridgelessness;
\item independently constructs its binary cycle space and all 900
      \(\mathrm{GL}(3,2)\)-orbits of nowhere-zero \(W\)-flows;
\item for each of the 21 flags, constructs the original vertex-potential
      affine system directly from local triangles and edge equations;
\item asserts that its scalar local affine rows recover the literal
      allowed potential set exactly;
\item separately constructs the component system of
      Theorem~\ref{thm:component}; and
\item checks both decisions, Lemma~\ref{lem:beta},
      \eqref{eq:dictionary}, and the switch update
      \eqref{eq:betaswitch} on all \(900\cdot21=18{,}900\) systems.
\end{enumerate}
The result is 18,900 agreements and 4,856 soluble systems.  This is a
finite regression audit, not a proof of the general theorem.

\subsection{Star-fibre selector census}

For a vertex-star packing \(P\), let \(f_P\) be its
fundamental-completion flow and let
\[
 \chi_{d,H}(P)=
 \one\!\left[\beta_H(P)\in
       \Span\{\gamma_C^d(P):C\in\C(P,H)\}\right].
\]
A possible universal selection target is to choose \(P\) with
\begin{equation}
       \bigoplus_{\substack{H\le W,\ \dim H=2\\0\ne d\in H}}
       \chi_{d,H}(P)=1.                                     \label{eq:selector}
\end{equation}
The xor in \eqref{eq:selector} has 21 terms and is
\(\mathrm{GL}(3,2)\)-invariant.  It implies that at least one parallel
system is soluble.  Universal existence of such a packing is open.

An AI-authored lazy-SAT producer exhaustively searched every feasible
rooted star fibre in every simple biconnected cubic graph through order
16.  The corpus has 4,461 graphs and 50,894 feasible rooted fibres; a
packing satisfying \eqref{eq:selector} was retained for every one.  The
producer uses exactly-one omission colourings, lazy connectivity cuts,
literal spanning-tree checks, fundamental completions, and binary affine
elimination.  This census is frozen separately at
\[
\texttt{../output/six-point-parallel-selector-through16/}.
\]
Its manifest and exact command are recorded in the accompanying
\texttt{README.md}.  The producer and the semantic checks are AI-authored,
and the exhaustive run does not have an independently checked SAT proof
certificate.  The result is therefore finite computational evidence, not
a theorem for arbitrary order and not independent human verification.

\section{Open selection problem}

Theorems~\ref{thm:component} and~\ref{thm:exact} isolate the remaining
universal problem without solving it.  For each plane \(H\), the defect
\(\beta_H\) depends only on the contracted \(E_H\)-component quotient,
while the three directions \(d\in H-\{0\}\) change the circuit columns.
One must select a flow---or, in the star-fibre approach, a packing---for
which at least one of the 21 defects lies in its corresponding span.

The parity selector \eqref{eq:selector} is deliberately stronger than
existence of one soluble flag.  Its finite success suggests searching
for a local exchange formula under basis exchanges of the three
spanning trees.  No such formula is proved here.  In particular, the
fixed-flow characterizations do not eliminate the existential quantifier
which carries the original difficulty of \FiveCDC.

\section*{Reproducibility and disclosure}

From this directory, run the fixed-flow checker and rebuild:
\begin{verbatim}
python3 verify.py --output RESULT.json
SOURCE_DATE_EPOCH=1785283200 tectonic main.tex
shasum -a 256 -c SHA256SUMS
\end{verbatim}
The selector package has its own frozen sources, corpus, result, census,
and manifest.  See \texttt{README.md} for its exact check command.

Every proof and program in this package was produced with material
assistance from OpenAI Codex agents using GPT-5-series models under
human direction.  ``Independent'' refers only to separation between
implementations and derivations inside the AI-assisted project.  It does
not mean independent human verification.  The literature comparison,
including the interpretation of Hu\v{s}ek--\v{S}\'amal
Theorem~3.16, was also AI-assisted and was checked against arXiv
version~1 dated July 27, 2026.  Specialist review remains necessary.

\bibliographystyle{plain}
\bibliography{references}

\end{document}
